\scanpage{229}
% consequence of (15)
(15) nous donne une \uline{réciproque} pour la condition envisagée pour que $(D_i,\unsure{\varphi_i|\omega_i})$ soit de Picard-Fuchs : il f. et s. que l'on ait (7). Cela résulte du th. de De Rham sur les fibres $X_s$ .....

% heading and computation with horizontal sections
\uline{Calcul de} $\mu(a,b)$ par voie ``algébrique'', sans calcul d'intégrale. \inserted{C'est local sur $S$.} La formule (5) nous ramène à calculer $\mu(a,b)$ \inserted{au voisinage de $s_0$,} pour \deleted{$a,b$} $a(s_0)$ et $b(s_0)$ assez voisins. Cela nous permet alors de supposer que sur $X/S$ il y a une connexion intégrable. On peut $\overline{\mathrm m}$ supposer, par un choix convenable, que les sections $a,b$ \unsure{soient|sont} horizontales. Alors on trouve (14)
\[
 \mu(a,b)=\varphi(b(s))-\varphi(a(s))
\]
où $\varphi$ est défini par (10).

% beginning of the nonhorizontal computation
Il y aurait lieu d'expliciter aussi une \unsure{forme|formule} explicite, sans supposer que $a(s)$ et $b(s)$ soient horizontales. Faisons le calcul au voisinage de $s_0$, et \unsure{soit|soient} $a_0,b_0$ les
% archive sheet number: 228
